%-------------------------------------------------------------------------------
%	SECTION TITLE
%-------------------------------------------------------------------------------
\cvsection{Expériences}

%-------------------------------------------------------------------------------
%	CONTENT
%-------------------------------------------------------------------------------
\begin{cventries}
    \cventry
    {Architecte et responsable technique$\sqbig$Transformation numérique} % Job title
    {Foyer} % Organization
    {Luxembourg} % Location
    {oct.\,2024 - présent} % Date(s)
    {
        \begin{cvitems} % Description(s) of tasks/responsibilities
            \item {Définition de l’architecture permettant la souscription et le paiement d’assurances en continu, malgré des systèmes historiques à disponibilité limitée, en combinant services redondés, traitements asynchrones et adaptation du parcours client.}
            \item {Planification des capacités et coordination technique pour sécuriser les échéances de livraison : estimation des charges, définition des compétences requises, réaffectation de ressources et évaluation technique des candidats.}
        \end{cvitems}
    }
%%---------------------------------------------------------
%

    \cventry
    {Architecte$\sqbig$ingénieur principal}
    {Disney}
    {Luxembourg}
    {oct.\,2022 -- sept.\,2024}
    {
        \begin{cvitems}
            \item {Conception et mise en œuvre d’évolutions des services d’identité portant leur capacité de 350 à 800 millions de comptes sans interruption, en coordination avec les équipes sécurité, mobiles et internet.}
            \item {Participation aux exercices mensuels de continuité régionale et coordination du transfert de connaissances d’un périmètre d’environ quinze personnes vers les équipes internes ; introduction de fonctions AWS Lambda en Rust réduisant d’environ 80 fois leur coût d’exécution unitaire mesuré.}
        \end{cvitems}
    }

%%---------------------------------------------------------
    \cventry
    {Architecte données}
    {Depop}
    {Luxembourg}
    {mars\,2022 -- sept.\,2022}
    {
        \begin{cvitems}
            \item {Définition d’un modèle harmonisé des clients, produits et ventes afin de réconcilier les données utilisées par l’application mobile, le service client et les équipes d’analyse.}
            \item {Conception et obtention de l’adhésion à une architecture en continu réduisant le délai de disponibilité des données de 24 heures à moins de 15 minutes, avec des transformations compréhensibles et contrôlables par les analystes.}
        \end{cvitems}
    }

%%---------------------------------------------------------

    \cventry
    {Ingénieur données$\sqbig$Architecte de solutions}
    {Zego}
    {Luxembourg}
    {nov.\,2021 -- août\,2022}
    {
        \begin{cvitems}
            \item {Intégration de données automobiles et assurantielles, et évaluation avec les équipes de science des données d’une trajectoire de modernisation des traitements Python sur Flink.}
        \end{cvitems}
    }
%
%%---------------------------------------------------------

    \cventry
    {Ingénieur données$\sqbig$Architecte de solutions}
    {Slice}
    {Londres}
    {avr.\,2021 -- oct.\,2021}
    {
        \begin{cvitems}
            \item {Définition d’une architecture Kafka Streams et ksqlDB réduisant la latence des données commerciales et la dépendance à des traitements Databricks coûteux.}
        \end{cvitems}
    }

%%---------------------------------------------------------
    \cventry
    {Architecte plateforme}
    {Compare the Market}
    {Londres}
    {janv.\,2020 -- mars\,2021}
    {
        \begin{cvitems}
            \item {Définition et adoption de standards d’ingénierie communs à deux équipes : automatisation des tests, développement piloté par les tests, spécifications comportementales et contrôles qualité, améliorant la fiabilité et la prévisibilité des livraisons.}
            \item {Pilotage technique de la consolidation des grappes Kafka et de leur migration vers AWS MSK, avec automatisation de l’infrastructure, contrôle des accès et traçabilité des données ; migration progressive réalisée sans interruption ni corruption.}
        \end{cvitems}
    }

%%---------------------------------------------------------
    \cventry
    {Architecte données}
    {Depop}
    {Londres}
    {juin\,2019 -- déc.\,2019}
    {
        \begin{cvitems}
            \item {Conception d’une architecture d’ingestion en continu depuis RabbitMQ vers Kafka et Redshift, réduisant le délai d’accès aux données d’activité de l’application mobile de 24 heures à environ 15 minutes.}
            \item {Présentation de la trajectoire au directeur technique et mise en place de tests automatisés, de contrôles de qualité, de rapports d’exécution et de mécanismes de traçabilité des données.}
        \end{cvitems}
    }

    \cventry
    {Architecte données$\sqbig$Ingénieur données}
    {The Guardian}
    {Londres}
    {oct.\,2018 -- mai\,2019}
    {
        \begin{cvitems}
            \item {Conception d’un lac de données et d’un catalogue AWS pour centraliser et rendre accessibles les données de revenus lecteurs et d’abonnements.}
        \end{cvitems}
    }

    \cventry
    {Architecte plateforme}
    {Sainsbury's}
    {Londres}
    {déc.\,2017 -- sept.\,2018}
    {
        \begin{cvitems}
            \item {Copilotage technique de la plateforme nationale de gestion des récompenses Nectar, traitant environ 70 Go par heure et réduisant le délai de traitement de 24 heures à moins de 4 secondes ainsi que les coûts d’environ 70\%.}
            \item {Conception d’une architecture de protection des données fondée sur la minimisation, le chiffrement, la pseudonymisation, la rétention et le contrôle des accès, validée par les responsables produit, juridique et sécurité.}
        \end{cvitems}
    }

    \cventry
    {Ingénieur principal Scala}
    {HMPO, IPT, \& VOA}
    {Londres}
    {nov.\,2015 -- nov.\,2017}
    {
        \begin{cvitems}
            \item {Contribution principale au développement du service national de renouvellement des passeports en ligne et conception de mécanismes sécurisés et résilients d’ingestion des données.}
        \end{cvitems}
    }

    \cventry
    {Ingénieur principal Scala}
    {Home Office (Border Force)}
    {Londres}
    {mars\,2014 -- oct.\,2015}
    {
        \begin{cvitems}
            \item {Direction technique et développement principal d’un système national de prévision des flux passagers, combinant horaires de vols, événements en temps réel et modèles opérationnels afin d’anticiper les effectifs nécessaires aux contrôles frontaliers.}
            \item {Déploiement dans les trois principaux aéroports britanniques après un mois de fonctionnement parallèle ; promotion comme système principal et absence de dépassement du seuil ministériel d’attente pendant les six mois suivants.}
        \end{cvitems}
    }

    \cventry
    {Ingénieur}
    {Autres}
    {Londres et Paris}
    {2002 -- 2014}
    {
        \begin{cvitems}
            \item {Différents postes d’ingénieur logiciel confirmé et de responsable technique auprès de Capco, News UK, HappyInc, Sky, DMC Digital, ITHR, TripAdvisor, SixAndCo, VoxMobili, Fullsix, BoxNewMedia, AVS Consulting.}
        \end{cvitems}
    }

\end{cventries}